\setcounter{equation}{0}
\section{History of Quantum Computation}


In 1982, acting independently, Richard Feynman and Paul Benioff has claimed that quantum computing could be possible. Simulating quantum systems on classical computers requires an exponentially large amount of time. Feynman's first motivation was these kind of operations, he proposed that a quantum system could simulate other quantum systems more efficiently \cite{spwc82}. Benioff's motivation was rather continuing the the trend of miniaturization in computer industry \cite{cps80}. As circuits continue to shrink they will eventually reach atomic level, and quantum mechanical effects would become a major limiting factor. So the only feasible solution was to construct a computer which would make use of quantum mechanical effects.

In 1985 David Deutsch published a paper \cite{qtct85} which proposed a quantum Turing machine, which made use of a before unnoticed property of quantum computers. Though severely underrated at its time, Deutsch's paper would later be considered a landmark in quantum computation theory. In his paper, Deutsch showed that it is possible to perform some tasks on quantum computers much more efficiently than classical computers. Deutsch's idea was to make use of a phenomenon which is known as quantum superposition. An isolated electron, if placed in a superposition, can travel along many different paths through a circuit. If the memory could be placed in a quantum superposition before applying the algorithm, it could be in multiple states simultaneously. So just one run-through the algorithm would calculate all computational paths simultaneously. This idea (Which was to be called \emph{quantum parallelism}) would be the foundation of all quantum algorithms.

The examples in Deutsch's paper were of limited use and due to their probabilistic nature, should be run several times to ensure the correct outcome. But these early hints in the field of quantum computation theory stirred the imaginations of many scientists, and about 10 years later Peter Shor of AT\&T labs came up with a factorization algorithm for quantum computers.

In 1994 Peter Shor of AT\&T developed a quantum computer algorithm for factorization \cite{ptafpf94} that is exponentially faster than any known classical algorithm. This discovery triggered a massive interest in quantum computation, because factorization is very hard to do classically, and all best known classical algorithms are super-polynomial. Factorization is known to be in the NP (Non-Deterministic Polynomial) complexity class, but it was believed that it was not in the P (Deterministic Polynomial) class and therefore thought to be intractably difficult on classical computers. These difficulties have been the basis of many modern crypto systems - including the widely used RSA (the Rivest-Shamir-Adleman protocol), which usually rely on the difficulty of factoring a very large number, which is the product of two prime numbers (secret keys). With Shor's algorithm, it has been shown that it is possible to factorize very large numbers with quantum computers, which would shake the security of known key exchange methods.

Only two years later Lov Grover of bell labs invented a quantum algorithm for fast searching of a database \cite{fqma96}. Although Grover's algorithm gives only a quadratic speed up over the best possible classical algorithm, the fact that searching forms the basis of many other algorithms, and the possibility to plugged in Grover's algorithm to existing algorithms to immediately increase their efficiency, increased the importance of Grover's discovery.

These algorithms show the capabilities of a quantum computer, and the main motivation behind the research into quantum computation. It is logical to expect even more powerful algorithms to be discovered in the near future. However, an implementation of a quantum computer which would be capable of convenient applications seems far away. Currently, quantum computers can operate only on a few qubits, and under extreme environmental conditions.




